%% Generated by tools/gen_error_codes.py from docs/errors/ of the compiler;
%% edit the compiler's pages or tools/error_themes.txt, then rerun it.

\clearpage
\section{Types and conversions}
\label{sec:codes:types}

\begin{longtable}{@{}l p{0.78\linewidth} r@{}}
\toprule Code & Message & Page\\ \midrule \endhead
\hyperref[sec:codes:e4049]{E4049} & array inner type cannot be \errhole{} & \pageref{sec:codes:e4049}\\
\hyperref[sec:codes:e4095]{E4095} & incompatible types \errhole{} and \errhole{} & \pageref{sec:codes:e4095}\\
\hyperref[sec:codes:e4096]{E4096} & incompatible values \errhole{} and \errhole{} & \pageref{sec:codes:e4096}\\
\hyperref[sec:codes:e4097]{E4097} & the type \errhole{} is not complete & \pageref{sec:codes:e4097}\\
\hyperref[sec:codes:e4217]{E4217} & type \errhole{} has no fields & \pageref{sec:codes:e4217}\\
\hyperref[sec:codes:e4218]{E4218} & temporary type \errhole{} has no size & \pageref{sec:codes:e4218}\\
\hyperref[sec:codes:e4220]{E4220} & type \errhole{} has no field named \errhole{} & \pageref{sec:codes:e4220}\\
\hyperref[sec:codes:e4227]{E4227} & cannot cast value of type \errhole{} into a value of type \errhole{} & \pageref{sec:codes:e4227}\\
\hyperref[sec:codes:e4270]{E4270} & expression produces a value, but is used as a type & \pageref{sec:codes:e4270}\\
\hyperref[sec:codes:e4272]{E4272} & expression produces a value of type \errhole{}, but is used as a type & \pageref{sec:codes:e4272}\\
\hyperref[sec:codes:e4273]{E4273} & expression produces a type, but is used as a value & \pageref{sec:codes:e4273}\\
\hyperref[sec:codes:e4275]{E4275} & expression produces the type \errhole{}, but is used as a value & \pageref{sec:codes:e4275}\\
\hyperref[sec:codes:e4278]{E4278} & void expression cannot be used as a value & \pageref{sec:codes:e4278}\\
\hyperref[sec:codes:e4304]{E4304} & \errhole{} has no type of its own, it takes the type of what it is assigned to & \pageref{sec:codes:e4304}\\
\bottomrule
\end{longtable}

\errorcode{E4049}{array inner type cannot be \errhole{}}
\label{sec:codes:e4049}

Raised during \textbf{validation}, from \texttt{ValidateErrorMessage::\allowbreak{}EMPTY\_\allowbreak{}ARRAY\_\allowbreak{}INNER} (\textit{src/\allowbreak{}ymirc/\allowbreak{}semantic/\allowbreak{}validator/\allowbreak{}errors.yr}).

\subsubsection*{What it means}

An array was declared with an inner type that cannot be stored in an array -- a type with no size, such as \tokennolst{void}.

\subsubsection*{Example}

\textit{test\_\allowbreak{}resources/\allowbreak{}diagnostics/\allowbreak{}test101.yr}:

%% check: skip
\begin{lstlisting}[style=coloredverbatimError]
// a static array of void values
fn nothing () {}

fn main () {
    let _ = [nothing (); 3];
}
\end{lstlisting}

\begin{lstlisting}[style=bashVerb, breaklines=true, escapechar=~]
Error[E4049] : array inner type cannot be void
 --> test_resources/diagnostics/test101.yr:(5,22)
 5  ~\lstbox{\textSFxi{}}~     let _ = [nothing (); 3];
    ~\lstbox{\textSFv{}}~                      ^
    ~\lstbox{\textSFxi{}}~
    = when validating test101::main -- (test_resources/diagnostics/test101.yr:(4,4))
    ~\lstbox{\textSFii{}}~~\lstbox{\textSFx{}}~~\lstbox{\textSFx{}}~~\lstbox{\textSFx{}}~~\lstbox{\textSFx{}}~~\lstbox{\textSFx{}}~~\lstbox{\textSFx{}}~~\lstbox{\textSFx{}}~
\end{lstlisting}

\subsubsection*{How to fix it}

Use a type with values. If the element type came from a call, check what the call returns.

\errorcode{E4095}{incompatible types \errhole{} and \errhole{}}
\label{sec:codes:e4095}

Raised during \textbf{validation}, from \texttt{ValidateErrorMessage::\allowbreak{}INCOMPATIBLE\_\allowbreak{}TYPE} (\textit{src/\allowbreak{}ymirc/\allowbreak{}semantic/\allowbreak{}validator/\allowbreak{}errors.yr}).

\subsubsection*{What it means}

Two types meet where they have to agree -- the two sides of an operator, a value and the variable it is assigned to, the branches of a conditional -- and neither converts to the other.

\subsubsection*{Example}

\textit{test\_\allowbreak{}resources/\allowbreak{}enumeration/\allowbreak{}test3.yr}:

%% check: skip
\begin{lstlisting}[style=coloredverbatimError]
enum X : f32
| A = 1;
\end{lstlisting}

\begin{lstlisting}[style=bashVerb, breaklines=true, escapechar=~]
Error[E4095] : incompatible types mut f32 and i32
 --> test_resources/enumeration/test3.yr:(2,1)
 2  ~\lstbox{\textSFxi{}}~ | A = 1;
    ~\lstbox{\textSFv{}}~ ^     ^
    ~\lstbox{\textSFii{}}~~\lstbox{\textSFx{}}~~\lstbox{\textSFx{}}~~\lstbox{\textSFx{}}~~\lstbox{\textSFx{}}~~\lstbox{\textSFx{}}~~\lstbox{\textSFx{}}~~\lstbox{\textSFx{}}~
\end{lstlisting}

\subsubsection*{How to fix it}

Convert one of them explicitly with \tokennolst{cast}, or change one of the two declarations. Mutability is part of the type, so two types differing only in it are also incompatible here.

\errorcode{E4096}{incompatible values \errhole{} and \errhole{}}
\label{sec:codes:e4096}

Raised during \textbf{validation}, from \texttt{ValidateErrorMessage::\allowbreak{}INCOMPATIBLE\_\allowbreak{}VALUES} (\textit{src/\allowbreak{}ymirc/\allowbreak{}semantic/\allowbreak{}validator/\allowbreak{}errors.yr}).

\subsubsection*{What it means}

Two values meet where they have to agree and cannot be reconciled, cf. \hyperref[sec:codes:e4095]{E4095}, the difference being in the values rather than in their types.

\subsubsection*{Example}

\textit{test\_\allowbreak{}resources/\allowbreak{}templates/\allowbreak{}implicit/\allowbreak{}test5.yr}:

%% check: skip
\begin{lstlisting}[style=coloredverbatimError]
fn foo {U, N : usize, J = N * 3} (a : [U ; N], z : [U ; J]) {
    a;
    z;
}

fn main () {
    foo ([1, 2], [1, 2, 3, 4]);
}
\end{lstlisting}

\begin{lstlisting}[style=bashVerb, breaklines=true, escapechar=~]
Error[E4096] : incompatible values 6us and 4us
 --> test_resources/templates/implicit/test5.yr:(2,29)
 2  ~\lstbox{\textSFxi{}}~ fn foo {U, N : usize, J = N * 3} (a : [U ; N], z : [U ; J]) {
    ~\lstbox{\textSFv{}}~                             ^
    ~\lstbox{\textSFii{}}~~\lstbox{\textSFx{}}~~\lstbox{\textSFx{}}~~\lstbox{\textSFx{}}~~\lstbox{\textSFx{}}~~\lstbox{\textSFx{}}~~\lstbox{\textSFx{}}~~\lstbox{\textSFx{}}~
\end{lstlisting}

\subsubsection*{How to fix it}

Make the two values agree.

\errorcode{E4097}{the type \errhole{} is not complete}
\label{sec:codes:e4097}

Raised during \textbf{validation}, from \texttt{ValidateErrorMessage::\allowbreak{}INCOMPLETE\_\allowbreak{}TYPE} (\textit{src/\allowbreak{}ymirc/\allowbreak{}semantic/\allowbreak{}validator/\allowbreak{}errors.yr}).

\subsubsection*{What it means}

A type is used where its size or its layout is needed, and it is not finished being validated.

\subsubsection*{Example}

\textit{test\_\allowbreak{}resources/\allowbreak{}lit\_\allowbreak{}class/\allowbreak{}ctors/\allowbreak{}test26.yr}:

%% check: skip
\begin{lstlisting}[style=coloredverbatimError]
fn foo (_ : &A)-> i32;
fn bar (_ : &A)-> i32;

class A {
    let x : i32;

    pub self ()
        with x = foo (self)
    {}
}

class B over A {
    let z : i32;

    pub self ()
        with z = bar (super)
    {}
}
\end{lstlisting}

\begin{lstlisting}[style=bashVerb, breaklines=true, escapechar=~]
Error[E4097] : the type &(test26::A)::fields is not complete
 --> test_resources/lit_class/ctors/test26.yr:(7,9)
 7  ~\lstbox{\textSFxi{}}~     pub self ()
    ~\lstbox{\textSFv{}}~         ^^^^  note: when creating the value self of &(test26::A)::fields
   ...
 --> test_resources/lit_class/ctors/test26.yr:(8,23)
 8  ~\lstbox{\textSFxi{}}~         with x = foo (self)
    ~\lstbox{\textSFv{}}~                       ^^^^
    ~\lstbox{\textSFii{}}~~\lstbox{\textSFx{}}~~\lstbox{\textSFx{}}~~\lstbox{\textSFx{}}~~\lstbox{\textSFx{}}~~\lstbox{\textSFx{}}~~\lstbox{\textSFx{}}~~\lstbox{\textSFx{}}~
\end{lstlisting}

\subsubsection*{How to fix it}

Break the cycle, cf. \hyperref[sec:codes:e4017]{E4017}.

\errorcode{E4217}{type \errhole{} has no fields}
\label{sec:codes:e4217}

Raised during \textbf{validation}, from \texttt{ValidateErrorMessage::\allowbreak{}TYPE\_\allowbreak{}HAS\_\allowbreak{}NO\_\allowbreak{}FIELDINFO} (\textit{src/\allowbreak{}ymirc/\allowbreak{}semantic/\allowbreak{}validator/\allowbreak{}errors.yr}).

\subsubsection*{What it means}

Field information was asked of a type that has no fields.

\subsubsection*{Example}

\textit{test\_\allowbreak{}resources/\allowbreak{}diagnostics/\allowbreak{}test72.yr}:

%% check: skip
\begin{lstlisting}[style=coloredverbatimError]
fn main () {
    let _ = __pragma!field_infos (i32);
}
\end{lstlisting}

\begin{lstlisting}[style=bashVerb, breaklines=true, escapechar=~]
Error[E4217] : type i32 has no fields
 --> test_resources/diagnostics/test72.yr:(2,13)
 2  ~\lstbox{\textSFxi{}}~     let _ = __pragma!field_infos (i32);
    ~\lstbox{\textSFv{}}~             ^^^^^^^^
    ~\lstbox{\textSFii{}}~~\lstbox{\textSFx{}}~~\lstbox{\textSFx{}}~~\lstbox{\textSFx{}}~~\lstbox{\textSFx{}}~~\lstbox{\textSFx{}}~~\lstbox{\textSFx{}}~~\lstbox{\textSFx{}}~
\end{lstlisting}

\subsubsection*{How to fix it}

Ask it of a record, an entity or a class.

\errorcode{E4218}{temporary type \errhole{} has no size}
\label{sec:codes:e4218}

Raised during \textbf{validation}, from \texttt{ValidateErrorMessage::\allowbreak{}TYPE\_\allowbreak{}HAS\_\allowbreak{}NO\_\allowbreak{}SIZE} (\textit{src/\allowbreak{}ymirc/\allowbreak{}semantic/\allowbreak{}validator/\allowbreak{}errors.yr}).

\subsubsection*{What it means}

The size of a temporary type was asked for, and it has none.

\subsubsection*{Example}

\textit{test\_\allowbreak{}resources/\allowbreak{}struct/\allowbreak{}test1.yr}:

%% check: skip
\begin{lstlisting}[style=coloredverbatimError]
in test1;

record A {
    let a : A;
}

record B {
    let c : C;
}

record C {
    let b : B;
}
\end{lstlisting}

\begin{lstlisting}[style=bashVerb, breaklines=true, escapechar=~]
Error[E4218] : temporary type error has no size
 --> test_resources/struct/test1.yr:(8,9)
 8  ~\lstbox{\textSFxi{}}~     let c : C;
    ~\lstbox{\textSFv{}}~         ^
    ~\lstbox{\textSFxi{}}~ Note : for the field c
    ~\lstbox{\textSFxi{}}~  --> test_resources/struct/test1.yr:(8,9)
    ~\lstbox{\textSFxi{}}~  8  ~\lstbox{\textSFxi{}}~     let c : C;
    ~\lstbox{\textSFxi{}}~     ~\lstbox{\textSFv{}}~         ^
    ~\lstbox{\textSFxi{}}~     ~\lstbox{\textSFxi{}}~ Note : within type test1::B
    ~\lstbox{\textSFxi{}}~     ~\lstbox{\textSFxi{}}~  --> test_resources/struct/test1.yr:(7,8)
    ~\lstbox{\textSFxi{}}~     ~\lstbox{\textSFxi{}}~  7  ~\lstbox{\textSFxi{}}~ record B {
    ~\lstbox{\textSFxi{}}~     ~\lstbox{\textSFxi{}}~     ~\lstbox{\textSFv{}}~        ^
    ~\lstbox{\textSFxi{}}~     ~\lstbox{\textSFxi{}}~     ~\lstbox{\textSFii{}}~~\lstbox{\textSFx{}}~~\lstbox{\textSFx{}}~~\lstbox{\textSFx{}}~~\lstbox{\textSFx{}}~~\lstbox{\textSFx{}}~~\lstbox{\textSFx{}}~~\lstbox{\textSFx{}}~
    ~\lstbox{\textSFxi{}}~     ~\lstbox{\textSFii{}}~~\lstbox{\textSFx{}}~~\lstbox{\textSFx{}}~~\lstbox{\textSFx{}}~~\lstbox{\textSFx{}}~~\lstbox{\textSFx{}}~~\lstbox{\textSFvii{}}~~\lstbox{\textSFx{}}~
    ~\lstbox{\textSFxi{}}~ Note : for the field b
    ~\lstbox{\textSFxi{}}~  --> test_resources/struct/test1.yr:(12,9)
    ~\lstbox{\textSFxi{}}~ 12  ~\lstbox{\textSFxi{}}~     let b : B;
    ~\lstbox{\textSFxi{}}~     ~\lstbox{\textSFv{}}~         ^
    ~\lstbox{\textSFxi{}}~     ~\lstbox{\textSFxi{}}~ Note : within type test1::C
    ~\lstbox{\textSFxi{}}~     ~\lstbox{\textSFxi{}}~  --> test_resources/struct/test1.yr:(11,8)
    ~\lstbox{\textSFxi{}}~     ~\lstbox{\textSFxi{}}~ 11  ~\lstbox{\textSFxi{}}~ record C {
    ~\lstbox{\textSFxi{}}~     ~\lstbox{\textSFxi{}}~     ~\lstbox{\textSFv{}}~        ^
    ~\lstbox{\textSFxi{}}~     ~\lstbox{\textSFxi{}}~     ~\lstbox{\textSFii{}}~~\lstbox{\textSFx{}}~~\lstbox{\textSFx{}}~~\lstbox{\textSFx{}}~~\lstbox{\textSFx{}}~~\lstbox{\textSFx{}}~~\lstbox{\textSFx{}}~~\lstbox{\textSFx{}}~
    ~\lstbox{\textSFxi{}}~     ~\lstbox{\textSFii{}}~~\lstbox{\textSFx{}}~~\lstbox{\textSFx{}}~~\lstbox{\textSFx{}}~~\lstbox{\textSFx{}}~~\lstbox{\textSFx{}}~~\lstbox{\textSFvii{}}~~\lstbox{\textSFx{}}~
    ~\lstbox{\textSFii{}}~~\lstbox{\textSFx{}}~~\lstbox{\textSFx{}}~~\lstbox{\textSFx{}}~~\lstbox{\textSFx{}}~~\lstbox{\textSFx{}}~~\lstbox{\textSFvii{}}~~\lstbox{\textSFx{}}~
\end{lstlisting}

\subsubsection*{How to fix it}

Ask the size of a concrete type instead.

\errorcode{E4220}{type \errhole{} has no field named \errhole{}}
\label{sec:codes:e4220}

Raised during \textbf{validation}, from \texttt{ValidateErrorMessage::\allowbreak{}TYPE\_\allowbreak{}NO\_\allowbreak{}FIELD} (\textit{src/\allowbreak{}ymirc/\allowbreak{}semantic/\allowbreak{}validator/\allowbreak{}errors.yr}).

\subsubsection*{What it means}

A field name was used on a type that has no field of that name, cf. \hyperref[sec:codes:e4019]{E4019}.

\subsubsection*{Example}

\textit{test\_\allowbreak{}resources/\allowbreak{}diagnostics/\allowbreak{}test73.yr}:

%% check: skip
\begin{lstlisting}[style=coloredverbatimError]
fn main () {
    let a = [1, 2, 3];
    let _ = a.nosuchfield;
}
\end{lstlisting}

\begin{lstlisting}[style=bashVerb, breaklines=true, escapechar=~]
Error[E4220] : type [i32 ; 3us] has no field named nosuchfield
 --> test_resources/diagnostics/test73.yr:(3,15)
 3  ~\lstbox{\textSFxi{}}~     let _ = a.nosuchfield;
    ~\lstbox{\textSFv{}}~               ^^^^^^^^^^^
    ~\lstbox{\textSFii{}}~~\lstbox{\textSFx{}}~~\lstbox{\textSFx{}}~~\lstbox{\textSFx{}}~~\lstbox{\textSFx{}}~~\lstbox{\textSFx{}}~~\lstbox{\textSFx{}}~~\lstbox{\textSFx{}}~
\end{lstlisting}

\subsubsection*{How to fix it}

Check the spelling and the type.

\errorcode{E4227}{cannot cast value of type \errhole{} into a value of type \errhole{}}
\label{sec:codes:e4227}

Raised during \textbf{validation}, from \texttt{ValidateErrorMessage::\allowbreak{}UNDEFINED\_\allowbreak{}CAST\_\allowbreak{}OP} (\textit{src/\allowbreak{}ymirc/\allowbreak{}semantic/\allowbreak{}validator/\allowbreak{}errors.yr}).

\subsubsection*{What it means}

A \tokennolst{cast} was asked for between two types with no conversion between them.

\subsubsection*{Example}

\textit{test\_\allowbreak{}resources/\allowbreak{}char\_\allowbreak{}type/\allowbreak{}test3.yr}:

%% check: skip
\begin{lstlisting}[style=coloredverbatimError]
in test3;

fn main () {
    let a = 'a'c8;
    let b = 'b'c16;
    let c = 'c'c32;

    cast!u16 (a);
    cast!u8 (b);
    cast!u8 (c);
    cast!u32 (a);

    cast!c16 (12u8);
    cast!c32 (12u16);
    cast!c8 (12u32);
}
\end{lstlisting}

\begin{lstlisting}[style=bashVerb, breaklines=true, escapechar=~]
Error[E4227] : cannot cast value of type c8 into a value of type u16
 --> test_resources/char_type/test3.yr:(8,5)
 8  ~\lstbox{\textSFxi{}}~     cast!u16 (a);
    ~\lstbox{\textSFv{}}~     ^^^^
    ~\lstbox{\textSFii{}}~~\lstbox{\textSFx{}}~~\lstbox{\textSFx{}}~~\lstbox{\textSFx{}}~~\lstbox{\textSFx{}}~~\lstbox{\textSFx{}}~~\lstbox{\textSFx{}}~~\lstbox{\textSFx{}}~
\end{lstlisting}

\subsubsection*{How to fix it}

Go through a type both convert to, or define the cast operator on the source type.

\errorcode{E4270}{expression produces a value, but is used as a type}
\label{sec:codes:e4270}

Raised during \textbf{validation}, from \texttt{ValidateErrorMessage::\allowbreak{}USE\_\allowbreak{}AS\_\allowbreak{}TYPE} (\textit{src/\allowbreak{}ymirc/\allowbreak{}semantic/\allowbreak{}validator/\allowbreak{}errors.yr}).

\subsubsection*{What it means}

An expression producing a value was written where a type is expected.

\subsubsection*{Example}

\textit{test\_\allowbreak{}resources/\allowbreak{}diagnostics/\allowbreak{}test82.yr}:

%% check: skip
\begin{lstlisting}[style=coloredverbatimError]
fn main () {
    let _ : 12 = 1;
}
\end{lstlisting}

\begin{lstlisting}[style=bashVerb, breaklines=true, escapechar=~]
Error[E4270] : expression produces a value, but is used as a type
 --> test_resources/diagnostics/test82.yr:(2,13)
 2  ~\lstbox{\textSFxi{}}~     let _ : 12 = 1;
    ~\lstbox{\textSFv{}}~             ^^
    ~\lstbox{\textSFii{}}~~\lstbox{\textSFx{}}~~\lstbox{\textSFx{}}~~\lstbox{\textSFx{}}~~\lstbox{\textSFx{}}~~\lstbox{\textSFx{}}~~\lstbox{\textSFx{}}~~\lstbox{\textSFx{}}~
\end{lstlisting}

\subsubsection*{How to fix it}

Write the type. When the two are spelled the same, the name resolves to the value; qualify it so it resolves to the type.

\errorcode{E4272}{expression produces a value of type \errhole{}, but is used as a type}
\label{sec:codes:e4272}

Raised during \textbf{validation}, from \texttt{ValidateErrorMessage::\allowbreak{}USE\_\allowbreak{}AS\_\allowbreak{}TYPE\_\allowbreak{}VAL} (\textit{src/\allowbreak{}ymirc/\allowbreak{}semantic/\allowbreak{}validator/\allowbreak{}errors.yr}).

\subsubsection*{What it means}

An expression produces a value of the reported type where a type is expected, cf. \hyperref[sec:codes:e4270]{E4270}.

\subsubsection*{Example}

\textit{test\_\allowbreak{}resources/\allowbreak{}diagnostics/\allowbreak{}test126.yr}:

%% check: skip
\begin{lstlisting}[style=coloredverbatimError]
// a function name used as a type
fn nothing () {}

fn main () {
    let _x_ : nothing = 1;
}
\end{lstlisting}

\begin{lstlisting}[style=bashVerb, breaklines=true, escapechar=~]
Error[E4272] : expression produces a value of type test126::nothing, but is used as a type
 --> test_resources/diagnostics/test126.yr:(5,15)
 5  ~\lstbox{\textSFxi{}}~     let _x_ : nothing = 1;
    ~\lstbox{\textSFv{}}~               ^^^^^^^
    ~\lstbox{\textSFxi{}}~
    = when validating test126::main -- (test_resources/diagnostics/test126.yr:(4,4))
    ~\lstbox{\textSFii{}}~~\lstbox{\textSFx{}}~~\lstbox{\textSFx{}}~~\lstbox{\textSFx{}}~~\lstbox{\textSFx{}}~~\lstbox{\textSFx{}}~~\lstbox{\textSFx{}}~~\lstbox{\textSFx{}}~
\end{lstlisting}

\subsubsection*{How to fix it}

Write the type itself.

\errorcode{E4273}{expression produces a type, but is used as a value}
\label{sec:codes:e4273}

Raised during \textbf{validation}, from \texttt{ValidateErrorMessage::\allowbreak{}USE\_\allowbreak{}AS\_\allowbreak{}VALUE} (\textit{src/\allowbreak{}ymirc/\allowbreak{}semantic/\allowbreak{}validator/\allowbreak{}errors.yr}).

\subsubsection*{What it means}

A type was written where a value is expected.

\subsubsection*{Example}

\textit{test\_\allowbreak{}resources/\allowbreak{}diagnostics/\allowbreak{}test127.yr}:

%% check: skip
\begin{lstlisting}[style=coloredverbatimError]
// a function pointer type used as a value
fn main () {
    let _ = fn (i32)-> i32;
}
\end{lstlisting}

\begin{lstlisting}[style=bashVerb, breaklines=true, escapechar=~]
Error[E4273] : expression produces a type, but is used as a value
 --> test_resources/diagnostics/test127.yr:(3,13)
 3  ~\lstbox{\textSFxi{}}~     let _ = fn (i32)-> i32;
    ~\lstbox{\textSFv{}}~             ^^
    ~\lstbox{\textSFxi{}}~
    = when validating test127::main -- (test_resources/diagnostics/test127.yr:(2,4))
    ~\lstbox{\textSFii{}}~~\lstbox{\textSFx{}}~~\lstbox{\textSFx{}}~~\lstbox{\textSFx{}}~~\lstbox{\textSFx{}}~~\lstbox{\textSFx{}}~~\lstbox{\textSFx{}}~~\lstbox{\textSFx{}}~
\end{lstlisting}

\subsubsection*{How to fix it}

Construct a value of that type, or use a value that already exists.

\errorcode{E4275}{expression produces the type \errhole{}, but is used as a value}
\label{sec:codes:e4275}

Raised during \textbf{validation}, from \texttt{ValidateErrorMessage::\allowbreak{}USE\_\allowbreak{}AS\_\allowbreak{}VALUE\_\allowbreak{}TYPE} (\textit{src/\allowbreak{}ymirc/\allowbreak{}semantic/\allowbreak{}validator/\allowbreak{}errors.yr}).

\subsubsection*{What it means}

An expression produces the reported type where a value is expected, cf. \hyperref[sec:codes:e4273]{E4273}.

\subsubsection*{Example}

\textit{test\_\allowbreak{}resources/\allowbreak{}diagnostics/\allowbreak{}test137.yr}:

%% check: skip
\begin{lstlisting}[style=coloredverbatimError]
// a trait used as a value
trait Shape {}

fn main () {
    let _ = Shape;
}
\end{lstlisting}

\begin{lstlisting}[style=bashVerb, breaklines=true, escapechar=~]
Error[E4275] : expression produces the type test137::Shape, but is used as a value
 --> test_resources/diagnostics/test137.yr:(5,13)
 5  ~\lstbox{\textSFxi{}}~     let _ = Shape;
    ~\lstbox{\textSFv{}}~             ^^^^^
    ~\lstbox{\textSFxi{}}~
    = when validating test137::main -- (test_resources/diagnostics/test137.yr:(4,4))
    ~\lstbox{\textSFii{}}~~\lstbox{\textSFx{}}~~\lstbox{\textSFx{}}~~\lstbox{\textSFx{}}~~\lstbox{\textSFx{}}~~\lstbox{\textSFx{}}~~\lstbox{\textSFx{}}~~\lstbox{\textSFx{}}~
\end{lstlisting}

\subsubsection*{How to fix it}

Construct a value of the type.

\errorcode{E4278}{void expression cannot be used as a value}
\label{sec:codes:e4278}

Raised during \textbf{validation}, from \texttt{ValidateErrorMessage::\allowbreak{}VOID\_\allowbreak{}VALUE} (\textit{src/\allowbreak{}ymirc/\allowbreak{}semantic/\allowbreak{}validator/\allowbreak{}errors.yr}).

\subsubsection*{What it means}

An expression producing nothing was used where a value is required.

\subsubsection*{Example}

\textit{test\_\allowbreak{}resources/\allowbreak{}diagnostics/\allowbreak{}test128.yr}:

%% check: skip
\begin{lstlisting}[style=coloredverbatimError]
// an array literal of void values
fn nothing () {}

fn main () {
    let _ = [nothing (), nothing ()];
}
\end{lstlisting}

\begin{lstlisting}[style=bashVerb, breaklines=true, escapechar=~]
Error[E4278] : void expression cannot be used as a value
 --> test_resources/diagnostics/test128.yr:(5,22)
 5  ~\lstbox{\textSFxi{}}~     let _ = [nothing (), nothing ()];
    ~\lstbox{\textSFv{}}~                      ^
    ~\lstbox{\textSFxi{}}~
    = when validating test128::main -- (test_resources/diagnostics/test128.yr:(4,4))
    ~\lstbox{\textSFii{}}~~\lstbox{\textSFx{}}~~\lstbox{\textSFx{}}~~\lstbox{\textSFx{}}~~\lstbox{\textSFx{}}~~\lstbox{\textSFx{}}~~\lstbox{\textSFx{}}~~\lstbox{\textSFx{}}~
\end{lstlisting}

\subsubsection*{How to fix it}

Use a call that returns something, or make the surrounding construct a statement.

\errorcode{E4304}{\errhole{} has no type of its own, it takes the type of what it is assigned to}
\label{sec:codes:e4304}

Raised during \textbf{validation}, from \texttt{ValidateErrorMessage::\allowbreak{}TYPEOF\_\allowbreak{}UNTYPED} (\textit{src/\allowbreak{}ymirc/\allowbreak{}semantic/\allowbreak{}validator/\allowbreak{}errors.yr}).

\subsubsection*{What it means}

\tokennolst{typeof} is applied to an untyped \tokennolst{null} or \tokennolst{none}. Neither literal has a complete type: \tokennolst{null} becomes the pointer type it is assigned to, and \tokennolst{none} the option type it is assigned to, so there is no type for \tokennolst{typeof} to name until then.

\subsubsection*{Example}

\textit{test\_\allowbreak{}resources/\allowbreak{}diagnostics/\allowbreak{}test148.yr}:

%% check: skip
\begin{lstlisting}[style=coloredverbatimError]
fn main () {
    let _ : typeof (null) = null;
}
\end{lstlisting}

\begin{lstlisting}[style=bashVerb, breaklines=true, escapechar=~]
Error[E4304] : null has no type of its own, it takes the type of what it is assigned to
 --> test_resources/diagnostics/test148.yr:(2,21)
 2  ~\lstbox{\textSFxi{}}~     let _ : typeof (null) = null;
    ~\lstbox{\textSFv{}}~                     ^^^^
    ~\lstbox{\textSFxi{}}~
    = when validating test148::main -- (test_resources/diagnostics/test148.yr:(1,4))
    ~\lstbox{\textSFii{}}~~\lstbox{\textSFx{}}~~\lstbox{\textSFx{}}~~\lstbox{\textSFx{}}~~\lstbox{\textSFx{}}~~\lstbox{\textSFx{}}~~\lstbox{\textSFx{}}~~\lstbox{\textSFx{}}~
\end{lstlisting}

\subsubsection*{How to fix it}

Name the type the literal is meant to have, or take the type of a value that already has it:

%% check: skip
\begin{lstlisting}[style=coloredverbatim]
let x : *i32 = null;
let _ : typeof (x) = null;
\end{lstlisting}

To recognize the \tokennolst{null} literal in a macro, match it with the \tokennolst{\_\_key} rule rather than comparing its type.
